\documentclass[11pt,a4paper]{article}
\usepackage[utf8]{inputenc}
\usepackage[T1]{fontenc}
\usepackage[catalan]{babel}
\usepackage{lmodern}
\usepackage[sfdefault,lining]{FiraSans}
\usepackage[a4paper,top=2.3cm,bottom=2.2cm,left=1.8cm,right=1.8cm,headheight=1.3cm,headsep=0.4cm,footskip=1cm]{geometry}
\usepackage{amsmath,amssymb}
\usepackage{icomma}
\usepackage{graphicx}
\usepackage[table]{xcolor}
\usepackage{tikz}
\usepackage{fontawesome5}
\usepackage{ccicons}
\usepackage{fancyhdr}
\usepackage{lastpage}
\usepackage{tasks}
\usepackage{needspace}
\usepackage[most]{tcolorbox}
\usepackage{array}
\usepackage[hidelinks,pdftitle={Fitxa d'exercicis de la unitat 7: Derivades i aplicacions. 1r de Batxillerat CCSS},pdfsubject={Matemàtiques Aplicades a les CCSS I},pdfkeywords={taxa de variació mitjana, derivada, funció derivada, regles de derivació, regla de la cadena, recta tangent, recta normal, monotonia, màxims i mínims, representació gràfica, optimització}]{hyperref}

\definecolor{navy}{HTML}{1B2748}
\definecolor{accent}{HTML}{E8603A}
\definecolor{teal}{HTML}{1F9E90}
\definecolor{soft}{HTML}{EEF1F8}
\definecolor{muted}{HTML}{5D6685}

\setlength{\parindent}{0pt}
\setlength{\parskip}{3pt}
\settasks{label=\alph*), label-format=\bfseries\color{navy}, label-width=1.3em, item-indent=1.8em, column-sep=0.6em, after-item-skip=5pt, before-skip=3pt, after-skip=4pt}

% ---------- capçalera i peu ----------
\newcommand{\llicencia}{\href{https://creativecommons.org/licenses/by-nc-sa/4.0/deed.ca}{\ccbyncsa}\ \ Llicència Creative Commons Reconeixement-NoComercial-CompartirIgual 4.0 Internacional (CC BY-NC-SA 4.0)}
\pagestyle{fancy}
\fancyhf{}
\fancyhead[L]{\raisebox{-0.35\height}{\includegraphics[height=0.95cm]{logo-gva.png}}}
\fancyhead[C]{\small\color{muted}Fitxa d'exercicis · Unitat 7: Derivades i aplicacions · 1r Batxillerat CCSS}
\fancyhead[R]{\raisebox{-0.35\height}{\includegraphics[height=0.95cm]{logo-vermellar.png}}}
\fancyfoot[L]{\scriptsize\color{muted}\llicencia}
\fancyfoot[R]{\small\color{muted}\thepage}
\renewcommand{\headrulewidth}{0.4pt}
\renewcommand{\headrule}{\color{muted!40}\hrule height \headrulewidth}
\fancypagestyle{primera}{\fancyhf{}\renewcommand{\headrulewidth}{0pt}
  \fancyfoot[L]{\scriptsize\color{muted}\llicencia}
  \fancyfoot[R]{\small\color{muted}\thepage}}

% ---------- nivells i exercicis ----------
\newcommand{\estrelles}[1]{{\color{accent}\ifcase#1\or\faStar\faStar[regular]\faStar[regular]\or\faStar\faStar\faStar[regular]\or\faStar\faStar\faStar\fi}}
\newcounter{exer}
\newcommand{\ex}[1]{\par\needspace{4\baselineskip}\medskip\refstepcounter{exer}\noindent{\color{navy}\bfseries\theexer.}\ \ }

\newcommand{\seccio}[2]{\par\needspace{16\baselineskip}\bigskip
  \begin{tcolorbox}[colback=navy,colframe=navy,coltext=white,arc=2mm,boxrule=0pt,left=3mm,right=3mm,top=1.2mm,bottom=1.2mm]
  {\large\bfseries #1\quad #2}
  \end{tcolorbox}\vspace{-2pt}}

\newcommand{\exf}[1]{\par\medskip\refstepcounter{exer}\noindent{\color{navy}\bfseries\theexer.}\ \ }
\providecommand{\gt}{>}\providecommand{\lt}{<}
% quadre de recordatori (només a la fitxa d'exercicis)
\newcommand{\recorda}[1]{\par\smallskip\begin{tcolorbox}[colback=accent!7,colframe=accent!55,arc=2mm,boxrule=0.5pt,left=2.5mm,right=2.5mm,top=1mm,bottom=1mm]\small{\color{accent}\faLightbulb[regular]}\ \textbf{Recorda:}\ #1\end{tcolorbox}}
\newcommand{\res}[1]{#1}\newcommand{\rest}[1]{#1}

\def\fmneed{11}% que el títol d'una secció no quede sol al final del full
% =====================================================================
% Capçalera FentMates de les fitxes: un color per a cada curs.
% Els colors són els de la pàgina inicial (LEVELS a assets/js/materials.js);
% el logo ve de _eines/logo/pdf/ (python3 _eines/logo/genera.py).
%
% Ús, al final del preàmbul de la unitat:
%   % =====================================================================
% Capçalera FentMates de les fitxes: un color per a cada curs.
% Els colors són els de la pàgina inicial (LEVELS a assets/js/materials.js);
% el logo ve de _eines/logo/pdf/ (python3 _eines/logo/genera.py).
%
% Ús, al final del preàmbul de la unitat:
%   % =====================================================================
% Capçalera FentMates de les fitxes: un color per a cada curs.
% Els colors són els de la pàgina inicial (LEVELS a assets/js/materials.js);
% el logo ve de _eines/logo/pdf/ (python3 _eines/logo/genera.py).
%
% Ús, al final del preàmbul de la unitat:
%   \input{../comu/fentmates.tex}
%   \fmcurs{1eso}
% i, al començament del document:
%   \capcalera{Matemàtiques}{1r ESO}{Fitxa d'exercicis · Unitat 2}{Divisibilitat}
% Opcional abans de \input: \def\fmneed{14} (espai mínim abans de cada secció).
% =====================================================================
\makeatletter
\def\fm@logodir{../../logo/pdf}
\providecommand{\fmneed}{10}
\def\fm@cursnom{}

% clau del curs, color i color del text damunt del color
\newcommand{\fm@curs}[3]{\@namedef{fm@color@#1}{#2}\@namedef{fm@text@#1}{#3}}
\fm@curs{1eso}{E6B84A}{navy}
\fm@curs{2eso}{3B82F6}{white}
\fm@curs{3eso}{FF7A50}{white}
\fm@curs{4eso-a}{A78BFA}{white}
\fm@curs{4eso-b}{6D28D9}{white}
\fm@curs{1batx-mi}{9B1C1C}{white}
\fm@curs{1batx-ccss}{EF6B6B}{white}
\fm@curs{2batx-mii}{1E6B47}{white}
\fm@curs{2batx-ccss}{2BB3A3}{white}

\newcommand{\fmcurs}[1]{%
  \@ifundefined{fm@color@#1}{\PackageError{fentmates}{Curs desconegut: #1}{}}{}%
  \edef\fm@tmp{\noexpand\definecolor{curs}{HTML}{\@nameuse{fm@color@#1}}%
               \noexpand\colorlet{curstext}{\@nameuse{fm@text@#1}}}\fm@tmp
  \colorlet{cursfosc}{curs!72!black}%
  \def\fm@logo{\fm@logodir/vertical-#1.pdf}%
  \def\fm@marca{\fm@logodir/marca-#1.pdf}}

% ---------- capçalera de la primera pàgina ----------
% #1 matèria (petit), #2 curs, #3 tipus de document i unitat, #4 títol
\newcommand{\capcalera}[4]{%
  \gdef\fm@cursnom{#2}%
  \par\noindent
  \begin{tcolorbox}[colback=curs,colframe=curs,boxrule=0pt,arc=2.2mm,width=3.5cm,height=2.15cm,
      valign=center,left=3mm,right=2mm,top=1mm,bottom=1mm,nobeforeafter,box align=center]
    \color{curstext}{\footnotesize #1}\par\vspace{1pt}{\fontsize{19}{21}\selectfont\bfseries #2\par}%
  \end{tcolorbox}\hspace{5mm}%
  \begin{minipage}[c]{\dimexpr\linewidth-3.5cm-5mm-2.9cm\relax}
    {\color{cursfosc}\small\bfseries\MakeUppercase{#3}}\par\vspace{2pt}
    {\color{navy}\fontsize{20}{23}\selectfont\bfseries #4\par}
  \end{minipage}\hfill
  \begin{minipage}[c]{2.6cm}\raggedleft\includegraphics[width=2.6cm]{\fm@logo}\end{minipage}\par
  \vspace{2pt}}

% ---------- barres de secció del color del curs ----------
\renewcommand{\seccio}[2]{\par\needspace{\fmneed\baselineskip}\bigskip
  \begin{tcolorbox}[colback=curs,colframe=curs,coltext=curstext,arc=2mm,boxrule=0pt,left=3mm,right=3mm,top=1.2mm,bottom=1.2mm]
  {\large\bfseries #1\quad #2}
  \end{tcolorbox}\vspace{-2pt}}

% ---------- estrelles de dificultat del color del curs (a Batxillerat no se'n posen) ----------
\@ifundefined{estrelles}{}{%
  \renewcommand{\estrelles}[1]{{\color{curs}\ifcase#1\or\faStar\faStar[regular]\faStar[regular]\or\faStar\faStar\faStar[regular]\or\faStar\faStar\faStar\fi}}}

% ---------- capçalera i peu de les pàgines ----------
% primera pàgina: logos institucionals al peu
\fancypagestyle{primera}{\fancyhf{}\renewcommand{\headrulewidth}{0pt}
  \fancyfoot[L]{\raisebox{-0.32\height}{\includegraphics[height=0.78cm]{logo-gva.png}}\hspace{5mm}%
    \raisebox{-0.32\height}{\includegraphics[height=0.78cm]{logo-vermellar.png}}}
  \fancyfoot[R]{\scriptsize\color{muted}\href{https://creativecommons.org/licenses/by-nc-sa/4.0/deed.ca}{\ccbyncsa}\ \ CC BY-NC-SA 4.0 · fentmates.org\qquad\small\thepage}}
% la resta: el símbol a l'esquerra i el curs a la dreta, amb el títol al centre com abans
\fancyhead[L]{\raisebox{-0.3\height}{\includegraphics[height=0.62cm]{\fm@marca}}}
\fancyhead[R]{\small\bfseries\color{cursfosc}\fm@cursnom}
\renewcommand{\headrule}{\color{curs!45}\hrule height \headrulewidth}
\makeatother

%   \fmcurs{1eso}
% i, al començament del document:
%   \capcalera{Matemàtiques}{1r ESO}{Fitxa d'exercicis · Unitat 2}{Divisibilitat}
% Opcional abans de \input: \def\fmneed{14} (espai mínim abans de cada secció).
% =====================================================================
\makeatletter
\def\fm@logodir{../../logo/pdf}
\providecommand{\fmneed}{10}
\def\fm@cursnom{}

% clau del curs, color i color del text damunt del color
\newcommand{\fm@curs}[3]{\@namedef{fm@color@#1}{#2}\@namedef{fm@text@#1}{#3}}
\fm@curs{1eso}{E6B84A}{navy}
\fm@curs{2eso}{3B82F6}{white}
\fm@curs{3eso}{FF7A50}{white}
\fm@curs{4eso-a}{A78BFA}{white}
\fm@curs{4eso-b}{6D28D9}{white}
\fm@curs{1batx-mi}{9B1C1C}{white}
\fm@curs{1batx-ccss}{EF6B6B}{white}
\fm@curs{2batx-mii}{1E6B47}{white}
\fm@curs{2batx-ccss}{2BB3A3}{white}

\newcommand{\fmcurs}[1]{%
  \@ifundefined{fm@color@#1}{\PackageError{fentmates}{Curs desconegut: #1}{}}{}%
  \edef\fm@tmp{\noexpand\definecolor{curs}{HTML}{\@nameuse{fm@color@#1}}%
               \noexpand\colorlet{curstext}{\@nameuse{fm@text@#1}}}\fm@tmp
  \colorlet{cursfosc}{curs!72!black}%
  \def\fm@logo{\fm@logodir/vertical-#1.pdf}%
  \def\fm@marca{\fm@logodir/marca-#1.pdf}}

% ---------- capçalera de la primera pàgina ----------
% #1 matèria (petit), #2 curs, #3 tipus de document i unitat, #4 títol
\newcommand{\capcalera}[4]{%
  \gdef\fm@cursnom{#2}%
  \par\noindent
  \begin{tcolorbox}[colback=curs,colframe=curs,boxrule=0pt,arc=2.2mm,width=3.5cm,height=2.15cm,
      valign=center,left=3mm,right=2mm,top=1mm,bottom=1mm,nobeforeafter,box align=center]
    \color{curstext}{\footnotesize #1}\par\vspace{1pt}{\fontsize{19}{21}\selectfont\bfseries #2\par}%
  \end{tcolorbox}\hspace{5mm}%
  \begin{minipage}[c]{\dimexpr\linewidth-3.5cm-5mm-2.9cm\relax}
    {\color{cursfosc}\small\bfseries\MakeUppercase{#3}}\par\vspace{2pt}
    {\color{navy}\fontsize{20}{23}\selectfont\bfseries #4\par}
  \end{minipage}\hfill
  \begin{minipage}[c]{2.6cm}\raggedleft\includegraphics[width=2.6cm]{\fm@logo}\end{minipage}\par
  \vspace{2pt}}

% ---------- barres de secció del color del curs ----------
\renewcommand{\seccio}[2]{\par\needspace{\fmneed\baselineskip}\bigskip
  \begin{tcolorbox}[colback=curs,colframe=curs,coltext=curstext,arc=2mm,boxrule=0pt,left=3mm,right=3mm,top=1.2mm,bottom=1.2mm]
  {\large\bfseries #1\quad #2}
  \end{tcolorbox}\vspace{-2pt}}

% ---------- estrelles de dificultat del color del curs (a Batxillerat no se'n posen) ----------
\@ifundefined{estrelles}{}{%
  \renewcommand{\estrelles}[1]{{\color{curs}\ifcase#1\or\faStar\faStar[regular]\faStar[regular]\or\faStar\faStar\faStar[regular]\or\faStar\faStar\faStar\fi}}}

% ---------- capçalera i peu de les pàgines ----------
% primera pàgina: logos institucionals al peu
\fancypagestyle{primera}{\fancyhf{}\renewcommand{\headrulewidth}{0pt}
  \fancyfoot[L]{\raisebox{-0.32\height}{\includegraphics[height=0.78cm]{logo-gva.png}}\hspace{5mm}%
    \raisebox{-0.32\height}{\includegraphics[height=0.78cm]{logo-vermellar.png}}}
  \fancyfoot[R]{\scriptsize\color{muted}\href{https://creativecommons.org/licenses/by-nc-sa/4.0/deed.ca}{\ccbyncsa}\ \ CC BY-NC-SA 4.0 · fentmates.org\qquad\small\thepage}}
% la resta: el símbol a l'esquerra i el curs a la dreta, amb el títol al centre com abans
\fancyhead[L]{\raisebox{-0.3\height}{\includegraphics[height=0.62cm]{\fm@marca}}}
\fancyhead[R]{\small\bfseries\color{cursfosc}\fm@cursnom}
\renewcommand{\headrule}{\color{curs!45}\hrule height \headrulewidth}
\makeatother

%   \fmcurs{1eso}
% i, al començament del document:
%   \capcalera{Matemàtiques}{1r ESO}{Fitxa d'exercicis · Unitat 2}{Divisibilitat}
% Opcional abans de \input: \def\fmneed{14} (espai mínim abans de cada secció).
% =====================================================================
\makeatletter
\def\fm@logodir{../../logo/pdf}
\providecommand{\fmneed}{10}
\def\fm@cursnom{}

% clau del curs, color i color del text damunt del color
\newcommand{\fm@curs}[3]{\@namedef{fm@color@#1}{#2}\@namedef{fm@text@#1}{#3}}
\fm@curs{1eso}{E6B84A}{navy}
\fm@curs{2eso}{3B82F6}{white}
\fm@curs{3eso}{FF7A50}{white}
\fm@curs{4eso-a}{A78BFA}{white}
\fm@curs{4eso-b}{6D28D9}{white}
\fm@curs{1batx-mi}{9B1C1C}{white}
\fm@curs{1batx-ccss}{EF6B6B}{white}
\fm@curs{2batx-mii}{1E6B47}{white}
\fm@curs{2batx-ccss}{2BB3A3}{white}

\newcommand{\fmcurs}[1]{%
  \@ifundefined{fm@color@#1}{\PackageError{fentmates}{Curs desconegut: #1}{}}{}%
  \edef\fm@tmp{\noexpand\definecolor{curs}{HTML}{\@nameuse{fm@color@#1}}%
               \noexpand\colorlet{curstext}{\@nameuse{fm@text@#1}}}\fm@tmp
  \colorlet{cursfosc}{curs!72!black}%
  \def\fm@logo{\fm@logodir/vertical-#1.pdf}%
  \def\fm@marca{\fm@logodir/marca-#1.pdf}}

% ---------- capçalera de la primera pàgina ----------
% #1 matèria (petit), #2 curs, #3 tipus de document i unitat, #4 títol
\newcommand{\capcalera}[4]{%
  \gdef\fm@cursnom{#2}%
  \par\noindent
  \begin{tcolorbox}[colback=curs,colframe=curs,boxrule=0pt,arc=2.2mm,width=3.5cm,height=2.15cm,
      valign=center,left=3mm,right=2mm,top=1mm,bottom=1mm,nobeforeafter,box align=center]
    \color{curstext}{\footnotesize #1}\par\vspace{1pt}{\fontsize{19}{21}\selectfont\bfseries #2\par}%
  \end{tcolorbox}\hspace{5mm}%
  \begin{minipage}[c]{\dimexpr\linewidth-3.5cm-5mm-2.9cm\relax}
    {\color{cursfosc}\small\bfseries\MakeUppercase{#3}}\par\vspace{2pt}
    {\color{navy}\fontsize{20}{23}\selectfont\bfseries #4\par}
  \end{minipage}\hfill
  \begin{minipage}[c]{2.6cm}\raggedleft\includegraphics[width=2.6cm]{\fm@logo}\end{minipage}\par
  \vspace{2pt}}

% ---------- barres de secció del color del curs ----------
\renewcommand{\seccio}[2]{\par\needspace{\fmneed\baselineskip}\bigskip
  \begin{tcolorbox}[colback=curs,colframe=curs,coltext=curstext,arc=2mm,boxrule=0pt,left=3mm,right=3mm,top=1.2mm,bottom=1.2mm]
  {\large\bfseries #1\quad #2}
  \end{tcolorbox}\vspace{-2pt}}

% ---------- estrelles de dificultat del color del curs (a Batxillerat no se'n posen) ----------
\@ifundefined{estrelles}{}{%
  \renewcommand{\estrelles}[1]{{\color{curs}\ifcase#1\or\faStar\faStar[regular]\faStar[regular]\or\faStar\faStar\faStar[regular]\or\faStar\faStar\faStar\fi}}}

% ---------- capçalera i peu de les pàgines ----------
% primera pàgina: logos institucionals al peu
\fancypagestyle{primera}{\fancyhf{}\renewcommand{\headrulewidth}{0pt}
  \fancyfoot[L]{\raisebox{-0.32\height}{\includegraphics[height=0.78cm]{logo-gva.png}}\hspace{5mm}%
    \raisebox{-0.32\height}{\includegraphics[height=0.78cm]{logo-vermellar.png}}}
  \fancyfoot[R]{\scriptsize\color{muted}\href{https://creativecommons.org/licenses/by-nc-sa/4.0/deed.ca}{\ccbyncsa}\ \ CC BY-NC-SA 4.0 · fentmates.org\qquad\small\thepage}}
% la resta: el símbol a l'esquerra i el curs a la dreta, amb el títol al centre com abans
\fancyhead[L]{\raisebox{-0.3\height}{\includegraphics[height=0.62cm]{\fm@marca}}}
\fancyhead[R]{\small\bfseries\color{cursfosc}\fm@cursnom}
\renewcommand{\headrule}{\color{curs!45}\hrule height \headrulewidth}
\makeatother

\fmcurs{1batx-ccss}

\begin{document}
\thispagestyle{primera}

\capcalera{Matemàtiques CCSS}{1r Batx}{Fitxa d'exercicis · Unitat 7}{\raggedright Derivades i aplicacions}
\vspace{10pt}
Nom i cognoms: \rule{8.4cm}{0.4pt}\hfill Grup: \rule{1.6cm}{0.4pt}\hfill Data: \rule{2.2cm}{0.4pt}
\vspace{8pt}

\begin{tcolorbox}[colback=soft,colframe=soft,arc=2mm,left=3mm,right=3mm,top=1.5mm,bottom=1.5mm]
\small Resol els exercicis al quadern mostrant tots els passos i simplifica sempre les derivades. Les seccions segueixen l'ordre dels apunts. Pots comprovar les derivades amb la calculadora de derivades dels apunts interactius.
\end{tcolorbox}

% =====================================================
\seccio{1}{Taxa de variació mitjana}

\exf{1} Troba la TVM de la funció $y=x^2-8x+12$ en els intervals $[1,2]$, $[1,5]$ i $[1,8]$. Després, troba-la en l'interval variable $[1,1+h]$ i comprova que, donant a $h$ els valors 1, 4 i 7, obtens els resultats anteriors.\recorda{$\text{TVM}[a,b]=\dfrac{f(b)-f(a)}{b-a}$ \ és el pendent del segment que uneix els punts $(a,f(a))$ i $(b,f(b))$. Si la TVM és positiva, la funció ha crescut de mitjana en l'interval.}

\ex{1} Llancem una pilota cap amunt i la seua altura (en metres) després de $t$ segons és $s(t)=20t-5t^2$. Calcula la velocitat mitjana en els intervals $[0,1]$, $[1,2]$ i $[2,4]$. Què indica el signe del darrer resultat?

\ex{1} Calcula la TVM de cada funció en l'interval indicat:\begin{tasks}(3)
  \task $f(x)=\sqrt x$ \ en $[1,4]$
  \task $g(x)=2^x$ \ en $[0,3]$
  \task $h(x)=\dfrac1x$ \ en $[1,2]$
\end{tasks}


% =====================================================
\seccio{2}{Derivada en un punt i funció derivada}

\exf{1} Calcula, aplicant la definició, les derivades següents:\begin{tasks}(1)
  \task La derivada de $y=\dfrac1x$ en el punt $x=-2$.
  \task La derivada de $y=-2x+4$ en els punts d'abscisses $-3$, $0$, $4$ i $7$. Per què ix sempre el mateix?
  \task La derivada de $y=3x^2-5x+1$ en els punts d'abscisses $-2$, $1$ i $0$.
\end{tasks}\recorda{$f'(a)=\displaystyle\lim_{h\to0}\frac{f(a+h)-f(a)}{h}$: és el pendent de la recta tangent a la gràfica en el punt $(a,f(a))$.}

\ex{2} Troba la funció derivada aplicant la definició $f'(x)=\displaystyle\lim_{h\to0}\frac{f(x+h)-f(x)}{h}$:\begin{tasks}(3)
  \task $f(x)=x^2+3x$
  \task $f(x)=\dfrac{2}{x+1}$
  \task $f(x)=\sqrt{x+1}$
\end{tasks}

\ex{1} Calcula per definició el pendent de la recta tangent a $y=x^2$ en $x=3$ i escriu-ne l'equació.


% =====================================================
\seccio{3}{Regles de derivació}

\exf{1} Calcula la funció derivada d'aquestes funcions simples:\begin{tasks}(2)
  \task $f(x)=2x^3-5x^2+3x-2$
  \task $f(x)=\sqrt{2x}+\sqrt[3]{5x^2}$
  \task $f(x)=\dfrac{1}{x\sqrt x}$
  \task $f(x)=\dfrac{x^3}{x^2+1}$
  \task $f(x)=x^2\cdot\ln x$
  \task $f(x)=\sqrt{3x^3}\cdot e^x$
\end{tasks}\recorda{$(x^n)'=nx^{n-1}$ (també amb exponents fraccionaris: escriu les arrels com a potències). \ $(f\cdot g)'=f'g+fg'$; \ $\left(\frac fg\right)'=\frac{f'g-fg'}{g^2}$.}

\ex{1} Deriva i simplifica:\begin{tasks}(3)
  \task $f(x)=x^4-\dfrac{3}{x^2}+5$
  \task $f(x)=\dfrac{x+1}{x-1}$
  \task $f(x)=\dfrac{x^2}{x+2}$
  \task $f(x)=\dfrac{e^x}{x}$
  \task $f(x)=\log_2x+3^x$
  \task $f(x)=(2x-1)\,e^x$
  \task $f(x)=(x^2+1)\,e^x$
  \task $f(x)=\dfrac{\ln x}{x}$
  \task $f(x)=\dfrac{x}{\ln x}$
\end{tasks}\recorda{$(e^x)'=e^x$, \ $(a^x)'=a^x\ln a$, \ $(\ln x)'=\frac1x$, \ $(\log_ax)'=\frac{1}{x\ln a}$.}


% =====================================================
\seccio{4}{Regla de la cadena}

\exf{1} Calcula la funció derivada d'aquestes funcions compostes:\begin{tasks}(2)
  \task $f(x)=\sqrt{x^2+1}$
  \task $f(x)=x\,e^{x^2}$
  \task $f(x)=\dfrac{2x}{(x^2+1)^2}$
  \task $f(x)=\ln(x^3-1)$
  \task $f(x)=\dfrac{x}{\sqrt{x^2+1}}$
  \task $f(x)=(2x+1)^4$
\end{tasks}\recorda{Regla de la cadena: $\left[f(g(x))\right]'=f'(g(x))\cdot g'(x)$. Derivem la funció «de fora» deixant igual la «de dins» i ho multipliquem per la derivada de la de dins.}

\ex{2} Deriva:\begin{tasks}(3)
  \task $f(x)=e^{3x^2-x}$
  \task $f(x)=(x^2-3x)^5$
  \task $f(x)=e^{\sqrt x}$
  \task $f(x)=\ln(e^x+1)$
  \task $f(x)=\dfrac{1}{(3x-2)^2}$
  \task $f(x)=\sqrt{\ln x}$
  \task $f(x)=\ln\dfrac{x-1}{x+1}$
  \task $f(x)=\log(x^2+1)$
  \task $f(x)=2^{x^2+1}$
\end{tasks}

\ex{2} Calcula $f'(1)$ si $f(x)=x^2e^{-x}$, i $g'(1)$ si $g(x)=\ln(x^2+3)$.


% =====================================================
\seccio{5}{Recta tangent i recta normal}

\exf{1} Calcula l'equació de la recta tangent i de la recta normal a les funcions següents en el punt indicat:\begin{tasks}(1)
  \task $f(x)=x^2+2x-1$ \ en el punt d'abscissa $x=1$
  \task $f(x)=2x^2-3x$ \ en el punt d'abscissa $x=-1$
  \task $f(x)=\dfrac{x+1}{x-1}$ \ en el punt d'abscissa $x=2$
\end{tasks}\recorda{Tangent: $y=f(a)+f'(a)(x-a)$. \ Normal (perpendicular a la tangent): $y=f(a)-\dfrac{1}{f'(a)}(x-a)$.}

\ex{1} Troba la recta tangent a $f(x)=\ln x$ en $x=1$ i la recta tangent a $g(x)=e^{2x}$ en $x=0$.

\ex{2} Donada $f(x)=x^3-3x^2+1$:\begin{tasks}(1)
  \task En quins punts la recta tangent és horitzontal?
  \task En quins punts la recta tangent és paral·lela a la recta $y=9x$? Escriu l'equació d'aquestes tangents.
\end{tasks}

\ex{2} Troba $a$ perquè la recta tangent a $f(x)=x^2+ax$ en $x=1$ tinga pendent 5. Escriu l'equació d'aquesta tangent.


% =====================================================
\seccio{6}{Monotonia i extrems}

\exf{1} Estudia la monotonia (intervals de creixement i de decreixement) i els extrems relatius de les funcions:\begin{tasks}(2)
  \task $f(x)=-x^2+5x-1$
  \task $f(x)=2x^3-3x^2-12x+5$
\end{tasks}\recorda{Calculem $f'(x)$, resolem $f'(x)=0$ i estudiem el signe de $f'$ en els intervals que determinen aquests punts i els que no són del domini: on $f'>0$ la funció creix i on $f'<0$ decreix. Si passa de créixer a decréixer hi ha un màxim relatiu; si passa de decréixer a créixer, un mínim.}

\ex{2} \label{ex:racional}Estudia la monotonia i els extrems relatius de:\begin{tasks}(3)
  \task $f(x)=\dfrac{x^2}{x-1}$
  \task $g(x)=x\,e^x$
  \task $h(x)=\ln(x^2+1)$
\end{tasks}

\ex{2} Troba els extrems relatius de $f(x)=x^4-8x^2+3$ amb el criteri de la segona derivada.

\ex{2} Calcula el màxim i el mínim absoluts de $f(x)=x^3-3x+2$ en l'interval $[-2,3]$.\recorda{Els extrems absoluts d'una funció contínua en $[a,b]$ són en els punts singulars de dins de l'interval o en els extrems $a$ i $b$: n'hi ha prou de comparar-ne els valors.}

\ex{3} Troba $a$ i $b$ perquè $f(x)=x^3+ax^2+b$ tinga un extrem relatiu en el punt $(2,\,-3)$. És un màxim o un mínim?


% =====================================================
\seccio{7}{Representació gràfica de funcions}

\exf{2} Estudia i representa la funció $f(x)=x^3-3x$: domini, punts de tall, simetria, asímptotes, monotonia i extrems.\recorda{Passos: domini, punts de tall, simetria, asímptotes, monotonia i punts singulars, signe de la funció, taula de valors i representació gràfica.}

\ex{2} Representa la funció $f(x)=\dfrac{x^2}{x-1}$ de l'exercici \ref{ex:racional}. Estudia'n abans les asímptotes.

\ex{2} Estudia i representa $f(x)=\dfrac{4x}{x^2+1}$.

\ex{2} Estudia i representa $f(x)=x^4-2x^2$.

\ex{2} Estudia i representa $f(x)=x^3-6x^2+9x$.

\ex{2} Estudia i representa $f(x)=\dfrac{x^2+4}{x}$.

\ex{2} Estudia i representa $f(x)=\dfrac{x^2}{x^2-4}$.

\ex{3} Estudia i representa $f(x)=\dfrac{(x-1)^2}{x+1}$, amb les asímptotes verticals i obliqües.

\ex{3} Estudia i representa $f(x)=\dfrac{x}{x^2-4}$. Té extrems relatius?


% =====================================================
\seccio{8}{Problemes d'optimització}

\exf{1} El benefici mensual d'una empresa, en milers d'euros, quan ven $x$ centenars d'unitats és $B(x)=-x^2+40x-300$. Estudia el creixement i el decreixement de $B$. Quantes unitats ha de vendre per obtindre el benefici màxim? Quin és aquest benefici? Per a quines vendes té pèrdues?\recorda{Si coneixem la funció: estudiem el signe de la derivada (on creix i on decreix), trobem els punts on $f'(x)=0$, justifiquem si són màxims o mínims i comprovem que la solució té sentit en el context (unitats positives, hores dins de l'interval\ldots).}

\ex{2} El cost total de produir $x$ unitats d'un article és $C(x)=x^2+100x+3\,600$ euros. Troba la funció de cost mitjà per unitat, $M(x)=\frac{C(x)}{x}$, estudia'n el creixement i el decreixement i calcula quantes unitats cal produir perquè siga mínim. Quin és aquest cost mitjà?

\ex{2} El nombre de persones connectades a una plataforma de vídeo (en milers) entre les 17 h i les 24 h és $N(t)=-2t^3+21t^2-36t+40$, on $t\in[0,7]$ són les hores que han passat des de les 17 h. Estudia quan augmenta i quan disminueix el nombre de connexions. A quina hora hi ha més persones connectades? I menys?

\ex{2} Les vendes d'un mòbil nou, en milers d'unitats al mes, $t$ mesos després de llançar-lo són $V(t)=\dfrac{40t}{t^2+16}$, $t\geq0$. Estudia el creixement i el decreixement de les vendes. En quin moment són màximes i quantes se'n venen? Què passa a llarg termini?

\ex{3} Una empresa estima que el benefici d'una campanya, en milers d'euros, si inverteix $x$ milers d'euros en publicitat, és $B(x)=10x\,e^{-0{,}2x}$, $x\geq0$. Estudia el creixement i el decreixement de $B$. Quant ha d'invertir per obtindre el benefici màxim? Quin és?

\ex{2} Una empresa ven cada unitat que fabrica a 120~€, i el cost de fabricar-ne $x$ al dia és $C(x)=2x^2+40x+300$ euros. Escriu la funció de benefici diari i calcula quantes unitats ha de fabricar per maximitzar-lo. Quin és el benefici màxim?\recorda{Si hem de construir la funció: \ ingressos = preu $\cdot$ quantitat venuda; \ benefici = ingressos $-$ costos. Tria bé la variable, escriu el preu i la quantitat en funció d'ella i indica quins valors pot prendre.}

\ex{2} Un hotel té 80 habitacions i, a 50~€ la nit, les té totes ocupades. Per cada euro que puja el preu, queda una habitació buida. Cada habitació ocupada li costa 10~€ de neteja i manteniment. Quin preu fa màxim el benefici? Quantes habitacions s'ocupen i quin és el benefici?

\ex{2} Un concert benèfic té 800 espectadors si l'entrada costa 20~€. Per cada euro que es rebaixa l'entrada, hi van 50 persones més. Quin preu fa màxima la recaptació? Quantes persones hi van i quant es recapta?

\ex{3} La demanda setmanal d'un producte és $q=600-20p$ unitats quan el preu és $p$ euros. Fabricar cada unitat costa 6~€ i, a més, l'empresa té uns costos fixos de 1\,000~€ setmanals. Escriu el benefici en funció del preu i troba el preu que el fa màxim. Quantes unitats es venen i quin és el benefici?

\ex{3} Una comerciant té 10\,000 kg de raïm que podria vendre hui a 2~€ el kg. El preu augmentarà 5 cèntims cada dia, però cada dia se'n fan malbé 10 kg per cada quadrat dels dies que transcorren ($10t^2$ kg). Si pot esperar com a molt 20 dies, quin dia li convé vendre?

\vfill
\begin{center}\color{muted}\small\faCheckCircle[regular]\ \ Ànim! Comprova les teues derivades amb la calculadora dels apunts interactius.\end{center}

\end{document}
